\section{Discussion}

The proposed method is designed to address the main limitation of the image-only H-CAST baseline in fine-grained bird classification. H-CAST provides a strong hierarchical visual backbone by organizing image regions from fine to coarse and producing predictions at species, family, and order levels \cite{park2025visually}. However, it does not explicitly use the attribute annotations available in CUB-200-2011 \cite{wah2011cub}. Since many bird species differ by subtle visual traits, this thesis introduces attribute-guided learning to encourage the species-level representation to encode discriminative bird characteristics.

\subsection{Rationale of the Proposed Design}

Species-level bird classification requires fine-grained local evidence. Two species from the same family may have similar global appearance, but differ in localized traits such as bill shape, wing pattern, crown color, breast markings, or tail shape. These traits are directly reflected in the CUB attribute annotations. Therefore, using attributes as auxiliary supervision is a natural way to guide the model toward visual cues that are relevant for species discrimination.

The proposed model attaches the attribute-guided module to the species-level branch because the species level is the finest and most difficult level in the hierarchy. Family and order classification are coarser tasks and may rely more on broader visual structures. In contrast, species recognition often depends on subtle part-level differences. For this reason, the attribute-guided fusion is applied to the species-level branch, while family and order predictions remain produced by the original H-CAST hierarchical stages.

The use of attribute-group queries is also motivated by the structure of the CUB attributes. Instead of predicting 312 independent binary attributes, the model operates on semantic attribute groups such as \textit{bill shape}, \textit{wing color}, \textit{head pattern}, and \textit{tail shape}. This group-wise formulation is more aligned with the organization of CUB attributes and allows each query to specialize in one semantic trait type. The group-specific attribute classifiers then provide auxiliary supervision for these attribute-oriented representations.

HSDAS further improves the attribute usage strategy by selecting attributes according to species-level sibling discrimination. Using all 312 attributes may introduce redundancy or weakly useful supervision. HSDAS reduces this issue by prioritizing attributes that distinguish a species from its sibling species, especially hard siblings with similar attribute profiles. This makes the selected attribute subset more aligned with the fine-grained nature of the classification task.

The attribute attention diversity loss is introduced to reduce redundancy among attribute-related attention patterns. Since different attribute groups may correspond to different bird parts or visual properties, they should ideally capture complementary evidence. The diversity regularizer encourages macro-part attention maps to cover different graph nodes, helping the model avoid relying on only one dominant region.

\subsection{Advantages}

The first advantage of the proposed method is that it does not require attribute labels during inference. Attribute annotations are used only during training to supervise the attribute branch. At test time, the model receives only an image and produces order, family, and species predictions. Therefore, the inference setting remains the same as the original image-only H-CAST model.

Second, the method is compatible with the official CUB-200-2011 train/test protocol. The model can be trained using the standard image labels and class-level attribute annotations provided by the dataset. No additional manual annotation, part localization label, or attribute input is required.

Third, the group-wise attribute design improves interpretability. Since each attribute query corresponds to a semantic group, the learned attention patterns can be analyzed according to bird traits such as head, wing, tail, and body-related properties. Although these attention maps are not guaranteed to be precise attribute localization maps, they provide useful qualitative evidence about how the model uses attribute-oriented information.

Fourth, HSDAS provides a more compact attribute subset than FullAttr-312. By selecting attributes that are useful for distinguishing sibling species, HSDAS is designed to reduce the attribute space while preserving discriminative semantic cues. This helps avoid using all attributes indiscriminately and provides a more task-specific attribute supervision mask.

Finally, the diversity loss encourages different attribute macro-parts to attend to complementary graph-node regions. This is useful because bird recognition often requires multiple cues from different parts of the body. Encouraging diversity among attribute attention maps is intended to reduce redundancy and make the attribute-guided representation more informative.

\subsection{Limitations}

Despite these advantages, the proposed method also has several limitations. First, it depends on the availability and quality of class-level attribute annotations. The approach is suitable for datasets such as CUB-200-2011, where attribute annotations are provided. However, it may not be directly applicable to datasets without reliable semantic attributes. Another limitation is that the attribute targets used in this thesis are class-level attributes. Therefore, all images from the same species share the same attribute target, even though the visibility of a trait may vary across images due to pose, occlusion, viewpoint, or image quality. This makes the attribute supervision suitable for injecting species-level semantic priors, but less capable of modeling image-specific attribute presence.

Second, HSDAS is designed for settings where both a taxonomy and class-level attributes are available. The method relies on sibling relationships in the taxonomy and compares class-level attribute profiles to identify hard siblings. Therefore, its usefulness may be limited in datasets without a meaningful hierarchy or without class-level semantic descriptions.

Third, the current attribute fusion is applied mainly to the species-level branch. This design is reasonable for fine-grained bird recognition, but it means that family and order classifiers do not directly receive fused attribute representations. Future work could study whether attribute guidance at coarser levels can further improve hierarchical consistency.

Fourth, the method does not use pixel-level attribute localization supervision. Although attribute-group queries attend to species-level graph nodes, the model is not explicitly trained to localize attributes in the image. Therefore, attention visualizations should be interpreted as internal attention patterns rather than ground-truth explanations of where each attribute appears.

Finally, the performance of the method may depend on the selected attribute mask and the loss weights. Different choices of \(\lambda_{attr}\), \(\lambda_{div}\), and the selected attribute subset can affect the balance between classification, attribute prediction, and attention diversity. This makes ablation studies necessary to understand the contribution of each component.

Overall, the proposed method extends H-CAST in a targeted way. It preserves the visually consistent hierarchical backbone, adds attribute-guided learning where fine-grained discrimination is most needed, and uses HSDAS to select a compact set of discriminative attributes. The method is therefore well suited to the CUB-200-2011 setting, where species recognition depends on subtle bird traits and where both taxonomy and attribute annotations are available.
